\section{Scene View dan Navigasi}

\subsection{Scene View}

Scene view memungkinkan Anda melihat dan mengedit dunia game secara visual. Tool navigasi (toolbar di atas Scene):

\begin{itemize}
  \item \textbf{Hand Tool (Q)}: Geser pandangan tanpa mengubah objek
  \item \textbf{Move (W)}: Pindahkan objek sepanjang sumbu X, Y, Z
  \item \textbf{Rotate (E)}: Putar objek
  \item \textbf{Scale (R)}: Ubah ukuran objek
  \item \textbf{Frame (F)}: Fokus kamera Scene pada objek terpilih
\end{itemize}

Navigasi orbit dan zoom dilakukan melalui kombinasi tombol Alt dengan klik mouse serta scroll.

\subsection{Konsep Project, Scene, dan Asset}

\textbf{Project} adalah kontainer untuk seluruh aset game (script, model, texture, material, scene). \textbf{Scene} adalah ruang kerja yang berisi GameObject dalam satu level atau layar. Project memiliki struktur folder (Scenes, Materials, Scripts, dll.) untuk mengorganisir aset. Scene disimpan sebagai file .unity di dalam project.

\subsection{Primitive 3D dan Material}

Unity menyediakan primitive 3D (Cube, Sphere, Plane, Cylinder, Capsule) sebagai geometri dasar. Setiap primitive memiliki Mesh Filter (data geometri) dan Mesh Renderer (menampilkan geometri ke layar). \textbf{Material} mendefinisikan tampilan permukaan (warna, tekstur); diterapkan ke Mesh Renderer. Material memisahkan tampilan dari geometri---satu material bisa dipakai banyak objek.

\subsection{Lighting}

\textbf{Directional Light} mensimulasikan sumber cahaya jauh (matahari). Properti Emission mengatur intensitas dan warna. Rotasi memengaruhi arah cahaya. Scene memerlukan minimal satu sumber cahaya agar objek 3D terlihat.
